% ============================================================================
% Chapter 32 --- External references and standards
% ----------------------------------------------------------------------------
% Purpose : annotated catalogue of normative references used throughout
%           the document. Closes Part IV.
% Page target: ~4 pages.
% Dependencies: cited throughout Parts II and III.
% ============================================================================

\chapter{External references and standards}
\label{ch:external-references}

This chapter is a curated index of the normative references on
which the architecture and the playbook rest. It is not a
literature review; the standards listed below are those that
appear in the contracts of Part~II or in the wave declarations of
Part~III. Each entry indicates briefly the role the reference plays
in this document; full citations are in the bibliography.

%-----------------------------------------------------------------------------
\section{NIST publications}
\label{sec:er-nist}

The NIST series provides several of the architectural anchors
this Reference draws on. The following are most directly used.

\paragraph{NIST CSF~2.0 --- Cybersecurity Framework~\cite{nist-csf}.} The five-function structure (Govern,
Identify, Protect, Detect, Respond, Recover) is referenced
implicitly throughout the architecture: Govern aligns with the
sovereignty grid and the playbook's wave-0 governance discipline;
Detect and Respond align with the observability and audit
chapters; Recover aligns with the reversibility playbook.

\paragraph{NIST SP 800-207 --- Zero Trust
Architecture~\cite{nist-zt}.} The zero-trust principles --- never
trust by location, verify continuously, assume breach --- align
with the TAC model's per-activity context and with the policy
layer's per-request decision. The SP~800-207 abstract decomposition
into PE (policy engine), PA (policy administrator), and PEP
(policy enforcement point) maps onto the architecture's policy
layer (\trchap{ch:policy}{}) and its enforcement points across the
execution plane.

\paragraph{NIST SP 800-53 Rev.~5~\cite{nist-sp800-53}.} The
control families (AC, SC, AU, IA, CM, IR among others) are
referenced as the source for several of the architecture's
contractual obligations: SC-7 (boundary protection) for
\trchap{ch:network}{}; AU-2 and AU-12 (audit events) for
\trchap{ch:observability}{}; IA-2 and IA-8 (identification and
authentication) for \trchap{ch:identity}{}.

\paragraph{NIST SP 800-162 --- Attribute-Based Access
Control~\cite{nist-sp800-162}.} The ABAC framework underpins the
policy layer's evaluation model: attributes from the subject
(identity), the resource (classification), and the environment
(zone, time, posture) compose the policy decision.

\paragraph{NIST SP 800-171 Rev.~3~\cite{nist-171r3}.} The CUI
(Controlled Unclassified Information) protection requirements
inform the data-classification discussion of
\trchap{ch:storage}{}, particularly for institutions with
United-States government research contracts.

\paragraph{NIST SP 800-160 Vol.~2~\cite{nist-sp800-160v2}.}
Cyber-resiliency techniques --- segmentation, contextual
awareness, dynamic positioning, deception --- are referenced as
the resilience underpinning of the architecture's
fault-tolerance posture.

\paragraph{NIST SP 800-215~\cite{nist-sp800-215}.} Trusted
infrastructure considerations relevant to the reference
architecture's hardware and platform choices.

\paragraph{NIST SP 1800 series.} The practice-guide series
(notably SP 1800-3 \emph{Attribute Based Access Control} and
SP 1800-19 \emph{Trusted Cloud}) provides worked-example
deployments for several of the patterns this Reference develops.
These guides are practitioner-oriented and complement the
SP 800 family.

%-----------------------------------------------------------------------------
\section{ISO/IEC publications}
\label{sec:er-iso}

\paragraph{ISO/IEC 27001:2022.} The information-security
management system standard frames the institutional context in
which the architecture operates. The control catalogue (Annex~A)
is a recurring reference for procurement clauses; the management
system clauses inform the wave-0 governance discipline.

\paragraph{ISO/IEC 27002.} The implementation guidance
accompanying ISO/IEC 27001 expands several controls into
operational practice; A.8.22 (network segregation) is referenced
in \trchap{ch:network}{} as the framework anchor for the
zone-based segmentation pattern.

\paragraph{ISO/IEC 27005.} Risk management for information
security; provides the methodology underpinning the wave-aware
risk register of \trchap{ch:cross-cutting}{}.

\paragraph{ISO/IEC 19790 and FIPS 140 series.} Cryptographic-module
security requirements frame the HSM choices of
\trchap{ch:crypto}{}, particularly for institutions with
regulatory obligations that name a specific FIPS level.

%-----------------------------------------------------------------------------
\section{European regulatory framework}
\label{sec:er-eu}

\paragraph{General Data Protection Regulation
(\gls{gdpr})~\cite{gdpr}.} The risk-based security obligation
(Art.~32), the data-protection-by-design obligation (Art.~25),
and the data-protection-impact-assessment obligation (Art.~35)
are the EU instance of the data-protection regime against which
the data dimension of the sovereignty grid is scored (generically,
\loc{data-protection-regime}).

\paragraph{NIS2 Directive~\cite{nis2}.} The cybersecurity
obligations on essential and important entities, particularly
the management-body accountability for security measures
(Art.~20), are the EU instance of the cybersecurity-baseline
regime (generically, \loc{cybersecurity-baseline-regime}) that
informs the institutional sovereignty dimension and the wave-0
governance discipline. The incident-reporting
obligations align with the audit and IR pipeline of
\trchap{ch:wave7-audit-ir}{}.

\paragraph{\emph{Schrems~II} judgement
(C-311/18)~\cite{schrems2020}.} The CJEU's clarification that
contractual safeguards alone do not satisfy data sovereignty
when data falls under a third-country surveillance regime is,
for adopters under the EU regime, the anchor of the data-dimension
scoring in the sovereignty grid; generically that scoring follows
\loc{cross-border-transfer-regime}.

\paragraph{EU AI Act (Regulation 2024/1689).} The risk-based
classification of AI systems intersects with the policy-plane
authorisation decisions for AI-related TACs (the
Researcher-LLM-API pilot scenario of \trchap{ch:wave5-pilot}{}
in particular).

\paragraph{eIDAS 2.} The renewed framework for electronic
identification and trust services informs the cryptographic
services chapter for institutions that provide EU-recognised
electronic credentials.

\paragraph{EDPB transfer
guidelines~\cite{edpb-guidelines}.} The European Data
Protection Board's recommendations on data transfers
operationalise the post-Schrems-II posture for the data-
dimension assessment.

%-----------------------------------------------------------------------------
\section{Sector publications}
\label{sec:er-sector}

\paragraph{EDUCAUSE~\cite{educause-top10}.} In North-American
higher education, the principal venue for IT discussion. Its Top 10 IT
issues, governance case material, and identity-and-access-
management practice papers inform several of the
institutional-precedent claims in
\trchap{ch:institutional-precedents}{}; equivalent sector bodies
(e.g.\ UCISA in the UK, regional consortia elsewhere) play the
analogous role in other jurisdictions.

\paragraph{REN-ISAC~\cite{ren-isac}.} In the North-American context,
the Research and Education Networks Information Sharing and Analysis
Center provides incident-response intelligence and coordination
patterns referenced in \trchap{ch:wave7-audit-ir}{}; the analogous
role elsewhere is the regional NREN CSIRT or national R\&E ISAC.

\paragraph{ENISA~\cite{enisa-risk}.} The European Union Agency
for Cybersecurity provides risk-management and incident-response
guidance; its publications underpin several institutional
sovereignty-dimension considerations.

\paragraph{Research-and-education-network community publications.}
The institution's national research and education network, and the
regional research-network community to which it belongs, publish
architecture reviews, operational case material, and federation
guidance (the applicable inter-institutional identity-federation
profile, e.g.\ eduGAIN~\cite{edugain}, the associated
collaboration-management service eduTEAMS~\cite{eduteams}, and the
community authentication-and-authorisation AARC
Blueprint~\cite{aarc-blueprint})
that inform the multiple references in this document to identity
federation, audit pipelines, and research-network patterns.

\paragraph{Verizon DBIR~\cite{verizon-dbir}.} The annual Data
Breach Investigations Report provides the empirical anchor for
several of the threat-pattern discussions in the policy and audit
chapters.

%-----------------------------------------------------------------------------
\section{Open-source specifications and projects}
\label{sec:er-oss}

The reference instantiation of Part~II can be realised with
open-source specifications and projects such as those listed
below. The list is illustrative rather than prescriptive or
exhaustive; it records representative projects of the kind
referenced in the Reference's contracts and instantiations, and
equivalent alternatives may be substituted on a like-for-like
basis. Project documentation
should be consulted for current capability statements; the
ecosystem evolves at a faster cadence than this Reference.

\paragraph{Identity layer.} Keycloak~\cite{keycloak} (federated
identity provider with OIDC and SAML2 support);
Shibboleth~\cite{shibboleth} (SAML2 federation, widely deployed in
academia); the applicable inter-institutional identity-federation
technical profile~\cite{edugain}
(inter-institutional federation policy framework).

\paragraph{Policy layer.} Open Policy Agent (OPA)~\cite{opa-project}
(CNCF graduated, Rego-based policy decision engine);
Cedar~\cite{cedar-project} (open formal-semantics policy language).

\paragraph{Network layer.} FreeRADIUS~\cite{freeradius} (open-source
RADIUS server); Open vSwitch~\cite{openvswitch} (OVS, programmable
switch); Cilium~\cite{cilium} (eBPF-based CNI for Kubernetes,
micro-segmentation); NETCONF/YANG~\cite{rfc-6241-netconf,rfc-7950-yang}
(IETF standards for declarative network configuration).

\paragraph{Compute and storage.} Kubernetes~\cite{kubernetes} (CNCF,
container orchestration); OpenStack~\cite{openstack} (OpenInfra
Foundation, IaaS); Ceph~\cite{ceph} (distributed storage);
Harbor~\cite{harbor} (CNCF, container registry);
SLSA~\cite{slsa-spec} (Supply-chain Levels for Software Artifacts,
OpenSSF); Sigstore~\cite{sigstore} (artifact signing and
verification); OCI image format~\cite{oci-image-spec}.

\paragraph{Endpoint.} Munki~\cite{munki} (macOS software
distribution); osquery~\cite{osquery} (cross-platform endpoint
posture queries); Ansible~\cite{ansible} (declarative configuration
management); MicroMDM~\cite{micromdm} / NanoMDM (open Apple MDM
servers).

\paragraph{Observability.} Vector~\cite{vector-project} (event
collection and routing); Apache Kafka~\cite{kafka} (durable event
ingestion); OpenSearch~\cite{opensearch} (analytical store); Grafana
Loki~\cite{grafana-loki} (log-line search);
Grafana~\cite{grafana} (dashboards);
OpenTelemetry~\cite{opentelemetry-spec} (CNCF, telemetry collection
standard); CloudEvents~\cite{cloudevents-spec} (CNCF, event format).

\paragraph{Cryptography.} HashiCorp Vault~\cite{vault} (secrets and
PKI); step-ca~\cite{step-ca} (Smallstep, ACME-compliant institutional
CA); SoftHSM~\cite{softhsm} (software HSM for non-production);
Nitrokey HSM (hardware HSM with PKCS\#11 support).

\paragraph{Orchestration.} Argo Workflows~\cite{argo-workflows} (CNCF,
workflow engine); Argo CD~\cite{argocd} / Flux (CNCF, GitOps
reconcilers); Crossplane~\cite{crossplane} (CNCF,
infrastructure-as-code beyond Kubernetes);
Temporal.io~\cite{temporal} (durable workflow engine);
Tekton~\cite{tekton} (CNCF, CI/CD pipelines).

\medskip

\noindent The standards and projects above illustrate one viable
backbone for the Reference's instantiation. The list is a working
artefact: projects gain or lose maturity over time, new projects
emerge, specifications are revised, and equivalent components may
be substituted. The phase-2 international university network would
be the natural locus for maintaining a current and collectively
curated version of this catalogue.
